\section{Prototype Mission Control Flow}

MANAR manages mission progression through deterministic control flow logic implemented in \texttt{mission.cpp}.

\subsection{State Representation \& Flag Logic}
Rather than a rigid finite state machine, the current C++ prototype tracks mission lifecycle transitions through explicit state flags:
\begin{itemize}
    \item \texttt{missionstarted} (\texttt{bool}): Indicates whether a mission is active.
    \item \texttt{searchplanlocked} (\texttt{bool}): Indicates whether search location coordinates are frozen.
    \item \texttt{enroute} (\texttt{bool}): Indicates active transit toward a target search sector.
    \item \texttt{destinationsearched} (\texttt{bool}): Set when lawnmower sweep over the active sector completes.
    \item \texttt{rescueefound} (\texttt{bool}): Flag representing target detection.
    \item \texttt{waitingforhelp} (\texttt{bool}): Enforces a stationary hold upon candidate detection.
    \item \texttt{isdestinationhome} (\texttt{bool}): Set when the active navigation target is the home base station.
\end{itemize}

\subsection{Sequential Sector Progression}
Multi-location search plans are executed sequentially. Each entry in \texttt{search\_locations.json} represents a search sector origin ($L_1, L_2, \dots, L_N$).

When lawnmower sweeps over location $k$ ($1 \le k \le N$) finish, \texttt{lawnmower()} sets \texttt{destinationsearched = true}. On the next loop iteration of \texttt{missionstatusupdater()}, the control core advances the index:
\begin{equation}
k \leftarrow k + 1
\end{equation}
resets lawnmower sweep variables (\texttt{searchrow = 0}, \texttt{horizontalmove = true}, \texttt{moveeast = true}), loads location $k+1$ as the new destination, and updates \texttt{runtime["search"]["current\_location\_id"] = k + 1}. Once all $N$ sectors have been searched ($k = N$), the system initiates Return-to-Launch navigation.

\subsection{Candidate Detection and Rescuee Representation}
In the current prototype, candidate victim detection is exercised using a manual \texttt{rescueefound} trigger flag. When set, \texttt{configurerescueestate()} executes \texttt{stopflight()}, halting horizontal movement, switching flight mode to \texttt{Stall}, turning off active payload components to conserve battery, and setting \texttt{waitingforhelp = true}.

Automated perception algorithms will replace this manual trigger flag when the ML detection pipeline is implemented.

\subsection{Return-to-Launch (RTL) Semantics}
MANAR's selected architecture defines Return-to-Launch (RTL) as an on-demand navigation command that sets or changes the target destination to the configured launch location. It is not conceptually a permanent returning mission mode.

\textit{Current Prototype Behavior:} In the current software implementation, \texttt{activateRTL()} handles RTL preemption by checking if the aircraft is landed (calling \texttt{launch()} and engaging configured active flight mode if grounded), clearing rescue flags (\texttt{rescueefound = false}, \texttt{waitingforhelp = false}), overwriting \texttt{search\_locations.json} with a single home destination entry (\texttt{planning\_mode = "rtl\_home"}), and locking \texttt{isdestinationhome = true}.
